\chapter{Il tessuto muscolare}

% ---------------------------------------------------------------------------
\section{Generalità}

I \textbf{tessuti} sono gruppi di cellule che lavorano in cooperazione e che
condividono caratteristiche morfologiche simili: la loro esistenza permette
una vera e propria \textbf{divisione del lavoro} fra le cellule
dell'organismo. Nel corpo umano si distinguono quattro tipi di tessuto:
\textbf{epiteliale}, \textbf{connettivo}, \textbf{muscolare} e
\textbf{nervoso}. Questo capitolo è dedicato al tessuto muscolare.

Il \textbf{muscolo} è l'organo responsabile della generazione di forza e
movimento: un muscolo scheletrico, ad esempio il bicipite brachiale, appare
rigonfio quando è \textit{attivo} (contratto) e rilasciato quando è
\textit{inactive} (a riposo). Ogni muscolo è costituito da un insieme di
\textbf{fibre muscolari} organizzate in fasci, avvolte da connettivo e
raggiunte da vasi sanguigni e da fibre nervose motorie; un prelievo bioptico
(ad esempio a livello del polpaccio) permette di campionare il tessuto per
l'analisi istologica. Al microscopio elettronico a scansione la sezione
trasversale di un muscolo rivela l'intreccio di fibre e connettivo che lo
compone, mentre al microscopio ottico le fibre osservate in sezione
longitudinale mostrano l'andamento parallelo caratteristico e la presenza di
nuclei allineati lungo il loro decorso.

% ---------------------------------------------------------------------------
\section{I tre tipi di tessuto muscolare}

Si distinguono tre tipi di tessuto muscolare: \textbf{scheletrico},
\textbf{cardiaco} e \textbf{liscio}.

\subsection{Tessuto muscolare scheletrico}

Le cellule del tessuto muscolare scheletrico sono \textbf{lunghe,
cilindriche, striate e multinucleate}. Si trovano in rapporto con il tessuto
connettivo e il tessuto nervoso, all'interno dei muscoli scheletrici come
quelli degli arti. Le loro funzioni sono molteplici: muovere o stabilizzare
la posizione dello scheletro; controllare l'ingresso e l'uscita
dell'apparato digerente, l'uscita dell'apparato respiratorio e di quello
urinario; generare calore; proteggere i visceri.

Le fibre muscolari scheletriche sono voluminose e prive di diramazioni;
contengono filamenti di actina e miosina disposti in modo da produrre le
\textbf{striature} caratteristiche, ben visibili al microscopio ottico. I
\textbf{nuclei}, numerosi per ciascuna fibra, sono disposti in posizione
periferica, appena al di sotto della membrana cellulare.

\subsection{Tessuto muscolare cardiaco}

Le cellule del tessuto muscolare cardiaco (\textbf{cardiomiociti})
differiscono da quelle scheletriche perché sono più piccole, presentano
diramazioni e possiedono un solo \textbf{nucleo}, posto centralmente; le
estremità della cellula appaiono spesso ramificate a Y. Le
\textbf{cellule muscolari cardiache} sono interconnesse tra loro dai
\textbf{dischi intercalari}, giunzioni specializzate che assicurano la
trasmissione meccanica ed elettrica tra una cellula e l'altra. Anche in
questo tessuto sono presenti filamenti di actina e miosina disposti in modo
da produrre le striature caratteristiche.

La sede di questo tessuto è esclusivamente il \textbf{cuore}, di cui
costituisce il miocardio; le sue funzioni sono spingere il sangue in circolo
e regolare la pressione idrostatica ematica.

\subsection{Tessuto muscolare liscio}

Le cellule del tessuto muscolare liscio sono \textbf{piccole, affusolate, non
striate}, con un singolo \textbf{nucleo} posto al centro della
\textbf{cellula muscolare liscia}; non possiedono diramazioni. Si trovano a
circondare le pareti dei vasi sanguigni e degli apparati digerente,
respiratorio, urinario e genitale. Le loro funzioni comprendono la
propulsione di cibo, urina e secrezioni, nonché la regolazione del calibro
delle vie respiratorie e dei vasi, contribuendo così a regolare il flusso
ematico tissutale.

% ---------------------------------------------------------------------------
\section{Le fibre muscolari lente}

Tra le fibre del muscolo scheletrico si distingue una popolazione di fibre
\textbf{lente}, caratterizzate da:

\begin{itemize}
  \item \textbf{metabolismo aerobico};
  \item \textbf{elevato contenuto di mioglobina};
  \item \textbf{numerosi mitocondri};
  \item \textbf{elevata vascolarizzazione};
  \item capacità di sostenere una \textbf{contrazione prolungata};
  \item colorazione rossastra, da cui la denominazione popolare di
    \textbf{«carni rosse»}.
\end{itemize}

Queste caratteristiche rendono le fibre lente particolarmente adatte
all'attività prolungata a bassa intensità, a differenza delle fibre veloci,
orientate invece verso contrazioni rapide e potenti ma poco resistenti alla
fatica.

% ---------------------------------------------------------------------------
\section{Le fasce corporee}

Le \textbf{fasce} sono formazioni di tessuto connettivo che conferiscono
forza e stabilità, mantengono gli organi interni nelle loro relative
posizioni e forniscono una via per la distribuzione dei vasi sanguigni e
linfatici e dei nervi. Procedendo dalla superficie corporea verso la cavità
interna, al di sotto della \textbf{cute} si distinguono tre livelli di fascia:

\begin{itemize}
  \item \textbf{Fascia superficiale}: si trova tra la cute e gli organi
    sottostanti; è composta da tessuto areolare e tessuto adiposo, ed è nota
    anche come strato sottocutaneo o \textbf{ipoderma}.
  \item \textbf{Fascia profonda}: forma una forte intelaiatura fibrosa
    interna, costituita da tessuto connettivo denso; è legata a capsule,
    tendini, legamenti e altre strutture di sostegno.
  \item \textbf{Fascia sottosierosa}: si trova tra le membrane sierose e la
    fascia profonda; è composta da tessuto areolare.
\end{itemize}

Procedendo ulteriormente in profondità, oltre la fascia sottosierosa si
incontra la \textbf{membrana sierosa} che riveste la \textbf{cavità
corporea}, mentre a livello della parete corporea la fascia profonda
avvolge anche le \textbf{coste}; la \textbf{membrana cutanea} costituisce
il rivestimento più esterno dell'intero sistema.

% ---------------------------------------------------------------------------
\section{Fisiologia della contrazione muscolare}

L'attivazione della fibra muscolare segue una sequenza di eventi ben
definita: uno \textbf{stimolo chimico o elettrico}, veicolato dalla giunzione
neuromuscolare (placca motrice) tra l'assone del motoneurone e la fibra
muscolare, genera una \textbf{risposta di membrana} sotto forma di
\textbf{potenziale d'azione}, che a sua volta innesca la \textbf{risposta
intracellulare}, ossia la \textbf{contrazione} della fibra.

Questo fenomeno può essere registrato sperimentalmente posizionando
elettrodi di superficie sul bicipite e sul tricipite (collegati a un
amplificatore biologico e a un capo di messa a terra), mentre il soggetto
compie una contrazione muscolare volontaria, ad esempio flettendo il
braccio in un gesto di posa da bodybuilder. Il segnale raccolto viene
visualizzato come un tracciato elettromiografico grezzo e come il suo
inviluppo integrato (\textit{integrated EMG}), che ne rappresenta
l'andamento nel tempo.

\subsection{Elettromiografia (EMG)}

L'\textbf{elettromiografia (EMG)} è una metodica neurofisiologica utilizzata
per studiare il sistema nervoso periferico dal punto di vista funzionale.
L'esame può essere condotto inserendo un elettrodo ad ago direttamente
all'interno del ventre muscolare (ad esempio a livello della gamba o
dell'avambraccio), registrando l'attività elettrica generata dalle unità
motorie durante la contrazione.

% ---------------------------------------------------------------------------
\section{Tecniche di imaging del tessuto muscolare}

\subsection{Risonanza magnetica (RM)}

La \textbf{risonanza magnetica} consente di visualizzare in dettaglio
l'anatomia dei singoli muscoli e la loro disposizione reciproca.

In una sezione assiale T1 della coscia sono identificabili, ad esempio: il
muscolo vasto laterale (1); i muscoli vasto mediale e vasto intermedio (2);
il muscolo retto femorale (3); il muscolo sartorio (4); il muscolo
adduttore lungo (5); il muscolo adduttore breve (6); il muscolo gracile (7);
il muscolo semitendinoso (8); il muscolo grande gluteo (9).

In una sezione coronale T1 della coscia sono invece riconoscibili: il
muscolo bicipite femorale (1); il muscolo vasto laterale (2); il muscolo
gemello inferiore (3); il muscolo otturatorio interno (4); il gruppo
adduttore (5); il muscolo gracile (6).

\subsection{Spettroscopia a risonanza magnetica (MRS)}

La \textbf{spettroscopia a risonanza magnetica del fosforo-31} (\textsuperscript{31}P
MRS) permette di studiare il metabolismo energetico del muscolo in condizioni
di riposo (\textit{at rest}). Lo spettro \textsuperscript{31}P registrato a
livello del polpaccio presenta picchi caratteristici corrispondenti al
\textbf{fosfato inorganico (Pi)}, ai \textbf{fosfodiesteri (PDE)}, alla
\textbf{fosfocreatina (PCr)} — il picco più intenso — e ai tre gruppi fosfato
dell'\textbf{ATP} (\(\gamma\), \(\alpha\), \(\beta\)).

L'acquisizione viene guidata da un'immagine di risonanza magnetica in
sezione trasversale del polpaccio, ottenuta a 1,5\,T, in cui viene
delimitato un riquadro corrispondente al volume di interesse da cui vengono
acquisiti i dati di spettroscopia protonica (\textsuperscript{1}H MRS) relativi ai lipidi
muscolari.

% ---------------------------------------------------------------------------
\section{Biopsia muscolare}

La \textbf{biopsia muscolare} è una procedura effettuata per valutare le
anomalie del sistema muscolo-scheletrico e per diagnosticare le malattie che
coinvolgono il tessuto muscolare. Il prelievo può essere eseguito con un
apposito dispositivo bioptico (ad esempio uno strumento a scatto tipo
"Magnum"), oppure per via chirurgica, isolando ed esponendo un piccolo
campo operatorio a livello del sito di prelievo mediante divaricatori.

\subsection{Rischi della procedura}

\begin{itemize}
  \item lividi e dolore al sito di biopsia;
  \item sanguinamento prolungato dal sito di biopsia;
  \item infezione del sito di biopsia.
\end{itemize}

\subsection{Esame microscopico del campione}

L'osservazione al microscopio ottico del campione bioptico permette di
valutare sia la morfologia longitudinale delle fibre, striate e allineate
parallelamente tra loro, sia la loro organizzazione in sezione trasversale:
in quest'ultima si riconoscono i profili poligonali delle singole fibre
muscolari, delimitati dal \textbf{perimisio (P)}, il tessuto connettivo che
avvolge i fascicoli muscolari, e attraversati da un piccolo vaso
\textbf{capillare (C)} che decorre nello spazio interfascicolare.

% ---------------------------------------------------------------------------
\section{Alterazioni patologiche del tessuto muscolare}

L'esame istologico del tessuto muscolare può rivelare segni di sofferenza o
di malattia, tra cui:

\begin{itemize}
  \item la \textbf{perdita di striature} nel muscolo, indice di
    disorganizzazione dell'apparato contrattile;
  \item un tessuto che mostra \textbf{degenerazione con infiltrati
    infiammatori}, segno di un processo patologico in atto a carico delle
    fibre muscolari.
\end{itemize}
